\documentclass{article}
\usepackage[margin=2cm]{geometry}
\usepackage{amsmath}
\usepackage{graphicx}
\usepackage{booktabs}
\usepackage[utf8]{inputenc}
\usepackage{ragged2e}
\setlength{\emergencystretch}{3em}
\usepackage{paracol}
\begin{document}
\sloppy

\noindent\hfill\includegraphics[height=1.8em]{assets/figure_header_logo.png}\par\noindent\hrule\vspace{4pt}

\begin{paracol}{2}
\RaggedRight

computing the parameters for $\boldsymbol { z } . { \mathbf { u } } ^ { \ast }$ Incorporating the autoregressive likelihood into yields:


\begin{equation}
\begin{array} { r } { p ( x _ { u } \mid x _ { o }, b ) = \left | \operatorname* { d e t } \frac { d q _ { z _ { u }, b } } { d x _ { u } } \right | \prod _ { i = 1 } ^ { | u | } p ( z _ { u } ^ { i } \mid z _ { u } ^ { i - 1 },..., z _ { u } ^ { 1 }, x _ { o }, b ), } \\ {... ~ }\end{array}
\end{equation}


where$| u |$is the cardinality of the unobserved covariates.


\subsection*{0.1Training and Best Guess Objective}


Since we follow a flow-based approach, we have access to.
the normalized conditional likelihoods; thus, we can train
our model by maximizing log likelihood. During training, if
we have access to complete training data, we will need to
manually create binary masks$b$based on some predefined
distribution$p . { k } \cdot p . { k }$$\mu . { \ell }$ $p . { k }$is typically chosen based on the application. For instance, Bernoulli random masks are commonly
used for real-valued vectors. Given binary masks, training.
data are divided into$x$ $x . { u }$and$x . { o }$and fed into the conditional model$\pmb { p } ( \pmb { x } . { \pmb { u } } \mid \pmb { x } . { \sigma }, \pmb { b } )$ If training data already contains
missing values, we can only train our model on the remaining covariates. As before, we manually split each data point
into two parts,  and$x . { u }$  based on a binary mask$x . { \alpha }$ $b$ Note
that dimensions in corresponding to the missing values are$^ h$
always set to . i.e., they are never observed during training..$^ 0$
In this setting, we will need another binary mask$. { t t }$indicating those dimensions that are not missing. Accordingly, we
define observed dimensions as$x . { a } ~ = ~ x \left [ b \right ]$ and unobserved
dimensions as${ \mathfrak { x } } . { u } = x \left [ m \odot ( 1b ) \right ]$ In this setting, we optimize$p ( \pmb { \tau } . { \pmb { u } } \mid \pmb { x } . { 0 }, \pmb { m }, b ).$ During testing, we set${ \mathfrak { b } } = { \mathfrak { m } }$that
is, the model imputes the missing values conditioned on all.
the remaining covariates.


Best Guess ObjectiveGiven a trained model, multiple imputations can be easily accomplished by drawing multiple
samples from the learned conditional distribution. However,
guess". Unfortunately, the analytical mean$\mathbb { E } . { \mathbf { p ( x . { u } | z . { u }, b ) } } | \pmb { x . { u } } |$
is not available for flow-based deep generative modeis. rurthermore, getting an accurate empirical estimate could be
prohibitive in high dimensional space. In this work, we
propose a robust solution that gives a single best guess
in terms of the MSE metric (it can be easily extended to
other metrics, e.g. an adversarial one). Specifically, we obtain our best guess by inverting the conditional transformation over the mean of the latent distribution, i.e., $q . { \pmb z . { o }, b } ^ { - 1 } ( \bar { z } ) =$
$q . { \pmb { x } . { \sigma }, b } ^ { - 1 } ( \mathbb { E } . { p ; \pmb { z } ( \pmb { z } | \mathbf { z } . { \sigma }, b ) } [ z ] ).$ The meanis analytically available$z$
tor Gaussian mixture base model. To ensure that this best
guess is close to unobserved values, we optimize with an
auxiliary MSE loss:


\begin{equation}
\dot { \mathcal { L } } = - \log ~ p ( x _ { u } \mid x _ { o }, b ) + \lambda \| q _ { x _ { o }, b } ^ { - 1 } ( \bar { z } ) - x _ { u } \| ^ { 2 }, \quad ( 2 )
\end{equation}


where$\lambda$ controls the relative importance of the auxiliary
objective. Note that we only penalize one particular point
from $\pmb { p } ( \pmb { x . { u } } \mid \pmb { x . { o } }, \pmb { b } )$to be close to$x . { \cdot u }$ Hence, it does not affect the diversity of the conditional distribution. Unlike our
arbritrary conditioning task, conditional generative modeling deals only with a fixed vector of external information
$y$ (such as a class label) to condition the likelihood of a

\switchcolumn
\RaggedRight

fixed set of covariates (such as pixels in an image). Typi-$x ^ { * }$
cally, approaches supplement the inputs of common generative models with some encoding of the conditional information. Conditional GANs , for example, extended GANs by
inputting a class label encoding into both the generator and
discriminator function, allowing these models to learn multiple conditional distributions. Similarly, a conditional form
of VAE was proposed by , which introduces an additional
network that learns the conditional prior${ \mathfrak { p e } } ( z \mid { \mathfrak { y } } )$ where z
is some latent encoding andis a class label.$\boldsymbol { H }$


\subsection*{0.2Arbitrary Conditioning Models}


Previous attempts to learn probability distributions conditioned on arbitrary subsets of known covariates include the
Universal Marginalizer , which is trained as a feed-forward
network to approximate the marginal posterior distribution
of each unobserved dimension conditioned on the observed
ones. During sampling, they propose to use a sequential
sampling mechanism by adding a newly sampled dimension
to the observed sets and running the network in an autoregressive manner. However, we observe that VAEAC suffers
from failure modes common in VAEs given that they optimize with respect to ELBO loss . In particular, VAEs sometimes sacrifice learning the true posterior for minimizing reconstruction loss, which often leads the model to generate
blurry samples for multimodal distributions. We extend the
RealNVP model by replacing all the coupling layers with
our proposed arbitrary conditional alternative. For the sake
of sampling efficiency, we use standard Gaussian base likelihoods here. Implementation details of and baselines are provided in Appendix . Figure. shows samples drawn from and
VAEAC. More samples are available in Appendix . We notice our model can generate coherent and diverse inpaintings for all three datasets and different masks. Compared
to VAEAC, our model is capable of generating sharp samples and restoring more details. Even when the missing rate
is high, can still generate decent inpaintings. To quantitatively evaluate our model, we report peak signal-to-noise
ratio (PSNR) and negative log-likelihoods (NLL) in Table..
We generate 5 different masks for each test image and report the average scores and the standard deviation. We note
that PSNR is a metric that prefers blurry images over sample diversity . Log-likelihood measures how well the model
matches the real conditional distribution, but may not correlate with visual quality . Hence, we evaluate the tradeoff between sample quality and diversity via the precision
and recall scores (PRD) . Since we cannot sample from the
groundtruth conditional distribution, we compute the PRD
score between the imputed joint distribution ${ \pmb { p } } [ { \pmb { x } } . { o } ] { \pmb { p } } ( { \pmb { x } } . { u } )$
${ \pmb x } . { \pmb { \sigma } } )$ and the true joint distribution$p ( { \pmb x } )$via 10,000 samples
trom them. The PRD scores for two distributions measure
how much of one distribution can be generated by another.
Higher recall means a greater portion of samples from the
true distribution${ \pmb { p } } ( { \pmb { x } } )$ are covered by$\pmb { p } ( \pmb { x . { o } } ) \pmb { p } ( \pmb { x . { u } } \mid \pmb { x . { o } } )$ and
similarly, higher precision means a greater poruon or samples from$p ( x . { o } ) p ( x . { o } \mid x . { o } )$ are covered by ${ \pmb { p } } ( { \pmb { x } } )$.We report
the$( F . { 8 }, F . { \frac { 1 } { 2 4 } } )$pairs in Table. to represent recall and precision,
respectively.

\end{paracol}

\end{document}
